\section{The matrix obstruction}\label{mr:section}

We now prove the conditional matrix statement.  Only the three
properties of the base graph enter this section; the Strong Arnold
Property is not needed.  Throughout, $v=ms$ denotes the order of the
clique blow-up.  Matrix traces and squared Frobenius norms will be divided
by $v$, even when the matrix under consideration has fewer than $v$
rows.

\begin{theorem}[Matrix obstruction]\label{thm:matrix-obstruction}
For every $a\in(0,1]$ there are constants $\beta_0(a)>0$ and
$m_0(a)\in\mathbb N$ with the following property.  Let
$0<\beta\le\beta_0(a)$, let $m\ge m_0(a)$, and let $G$ be an
$m$-vertex graph satisfying the following conditions:
\begin{enumerate}[label=\textup{(\roman*)}]
\item Every matching of size at least $m/20$, partitioned into classes
of size at most $am$, contains two conflicting edges in different
classes.
\item Every matching of size at least $\beta m$ and every vertex set of
size at least $\beta m$ contain a conflicting edge--point pair.
\item Every two vertex sets of size at least $\beta m$ have a cross hole.
\end{enumerate}
Then there is $s_0=s_0(G,a,\beta)$ such that, for every integer
$s\ge s_0$, every real symmetric well-signed matrix $M$ for the clique
blow-up $G[K_s]$ with exactly one negative eigenvalue satisfies
\[
  \rank M>\frac{ms}{2}+1.
\]
Consequently $\mu(G[K_s])+1<ms/2$ for all these $s$.
\end{theorem}

The proof uses two normalizations.  A first normalization, separately
inside the vertex cliques, produces almost disjoint coordinate pairs.
A second normalization, inside groups of pairs having the same two base
labels, supplies additional positive directions.  The tree lemmas ensure
that both normalizations can be used in one bounded Gram matrix.  We first
give the linear-algebra estimates that make the errors independent of the
bounds furnished by the tree.

\subsection{Positive blocks and bounded spectral limits}

For a real symmetric matrix $A$, write $n_+(A)$ and $n_-(A)$ for its
positive and negative indices, counting multiplicity.  The notation
$\|A\|$ means the Euclidean operator norm, and
$\|A\|_{\mathrm F}^2=\Tr(A^{\mathsf T}A)$.  A symmetric matrix with
at most one negative eigenvalue is a Gram matrix for the Lorentz form
\[
  \langle x,y\rangle=-x_0y_0+\sum_{i\ge1}x_i y_i:
\]
diagonalize the matrix and multiply each eigenvector coordinate by the
square root of the absolute value of its eigenvalue.  Unused positive
coordinates can be added.  Every restriction or coefficient pullback of
this Gram matrix still has at most one negative eigenvalue.  For lists
$X,Y$ of Lorentz vectors, $\langle X,Y\rangle$ denotes their rectangular
Gram matrix; the same letter denotes the corresponding synthesis map.

The next facts are standard consequences of symmetric $M$-matrix
theory and the Perron--Frobenius theorem; see
\citet[Theorem~5.2]{FiedlerSchneider1983Analytic} and
\citet[Theorem~I]{DebreuHerstein1953Nonnegative}.
Here an $M$-matrix has the form $cI-A$ with $A$ entrywise
nonnegative and $c$ at least its spectral radius. We include the
elementary proofs for the precise symmetric cases used below.

\begin{lemma}[Normalization of a positive block]\label{mr:positive-block}
Let $D$ be real symmetric with nonpositive off-diagonal entries.
\begin{enumerate}[label=\textup{(\roman*)}]
\item If $D$ is positive definite, its positive inverse square root
$D^{-1/2}$ is entrywise nonnegative and has positive diagonal entries.
\item Suppose $D$ is positive semidefinite and the graph of its strictly
negative off-diagonal entries is connected.  If $D$ is singular, its
kernel is one-dimensional and is spanned by a vector all of whose
coordinates are positive.  Every proper principal submatrix of $D$ is
positive definite.
\end{enumerate}
\end{lemma}

\begin{proof}
For (i), choose $c>\lambda_{\max}(D)$ and put $A=I-D/c$.
All entries of $A$ are nonnegative, and its eigenvalues lie in $[0,1)$.
The convergent binomial series gives
\begin{equation}\label{mr:inverse-root}
 D^{-1/2}=c^{-1/2}\sum_{k=0}^{\infty}
       \frac{\binom{2k}{k}}{4^k}A^k.
\end{equation}
Every term is entrywise nonnegative, and the constant term has positive
diagonal.

For (ii), take $c>\lambda_{\max}(D)$ and consider $A=cI-D$.
It is symmetric and entrywise nonnegative, with positive diagonal and
connected off-diagonal support.  A unit vector maximizing its Rayleigh
quotient may be replaced by its coordinatewise absolute value, so a
nonnegative maximizing eigenvector exists.  Its eigenvalue is positive.
The eigenvector equation and connectivity imply that it has no zero
coordinate: a zero coordinate would force zeros at every adjacent
coordinate and then everywhere.  Moreover, every vector for the largest
eigenvalue has one sign.  Indeed, equality in the absolute-value
Rayleigh inequality forces equal signs across each positive edge, while
its absolute value is a nonnegative maximizing eigenvector and thus has
full support.  If the largest eigenspace had dimension at least two, it
would contain a nonzero vector orthogonal to the positive eigenvector;
that vector could not have one sign.  The largest eigenvalue is therefore
simple.

When $D$ is singular and positive semidefinite, this largest eigenvalue
of $A$ is $c$, so the asserted kernel description follows.  If a proper
principal submatrix were singular, extend a nonzero vector of zero
quadratic form by zero coordinates.  Its quadratic form under $D$ is
zero, hence it belongs to $\ker D$ by the spectral theorem.  This
contradicts full support.
\end{proof}

\begin{lemma}[Spectral support from signs and inertia]
\label{mr:spectral-support}
Let $v_j\to\infty$, and let $B_j$ be symmetric entrywise nonnegative
matrices of orders $d_j\le v_j$, with zero diagonal and
$\sup_j\|B_j\|<\infty$.  Suppose $n_-(I-B_j)\le1$.
Every weak subsequential limit of
\[
 \nu_j=\frac1{v_j}\sum_{i=1}^{d_j}\delta_{\lambda_i(B_j)}
\]
is supported on $[-1,1]$ and has first moment zero.  In particular, for
each $\varepsilon>0$, the number of eigenvalues of $I-B_j$ outside
$[-\varepsilon,2+\varepsilon]$ is $o(v_j)$.

If also
\[
 \liminf_j\frac{d_j}{v_j}\ge1-\delta,
 \qquad \rank(I-B_j)\le v_j/2+1,
\]
then
\begin{equation}\label{mr:spectral-error}
 \limsup_j\frac{\|B_j^2-I\|_{\mathrm F}^2}{v_j}\le2\delta.
\end{equation}
The constant in this inequality is independent of the operator bound.
\end{lemma}

\begin{proof}
By \Cref{lem:compact-measures}, all measures have bounded mass and support in a common compact interval,
so every subsequence has a weakly convergent further subsequence.  Fix
such a limit $\nu$, and write $\theta=\nu(\mathbb R)\le1$.
At most one eigenvalue of $B_j$ exceeds $1$, hence $\nu$ is supported
on $(-\infty,1]$.  For every integer $k\ge0$, entrywise nonnegativity
gives $\Tr B_j^{2k+1}\ge0$, and therefore
$\int x^{2k+1}\,d\nu(x)\ge0$.  If $\nu$ gave mass $h>0$ to
$(-\infty,-1-\varepsilon]$, this integral would be at most
$\theta-h(1+\varepsilon)^{2k+1}$, which is negative for large $k$.
This proves the support assertion.  The zero diagonal gives
$\int x\,d\nu(x)=0$.  The asserted eigenvalue count follows by applying
these facts to any putative subsequence on which the normalized
exceptional count stays positive.

For the additional assertion, pass to a subsequence realizing the
limsup in \eqref{mr:spectral-error} and then to a weak limit as above.
The nullity assumption supplies at least $d_j-v_j/2-1$ eigenvalues
equal to $1$.  Since singleton sets are closed, weak convergence gives
$r:=\nu(\{1\})\ge\theta-1/2$.  Consequently
\[
 \int_{x\ne1}(1+x)\,d\nu(x)=\theta-2r
       \le1-\theta\le\delta.
\]
For $-1\le x\le1$ we have
$(1-x^2)^2\le1-x^2\le2(1+x)$, and the first expression vanishes at
$x=1$.  Polynomial integrals converge on the common compact support,
so
\[
 \lim_j\frac{\|B_j^2-I\|_{\mathrm F}^2}{v_j}
  =\int(1-x^2)^2\,d\nu(x)\le2\delta.
\]
In this argument the bound on $\|B_j\|$ is fixed before $j$ tends to
infinity.  It removes the contribution of any exceptional eigenvalue
to each fixed moment; it does not enter the resulting bound.
\end{proof}

Row-maximum rounding also appears in the study of nonnegative
orthogonality constraints \citep{JiangMengWenChen2023Nonnegative,ChenHeZhang2025Stiefel}.
Here symmetry and zero diagonal produce a partial matching, with
bounds independent of the matrix order after normalization by $v$.

\begin{lemma}[Extraction of coordinate pairs]\label{mr:matching}
Let $B$ be a symmetric entrywise nonnegative $d\times d$ matrix with
zero diagonal, where $d\le v$.  Suppose
\[
 \|B^2-I\|_{\mathrm F}^2\le\varepsilon v,
 \qquad 0<\varepsilon\le 8^{-4},
 \qquad \tau=\varepsilon^{1/4}.
\]
There is a matching of coordinates with adjacency matrix $W$, extended
by zero on the omitted coordinates, such that
\begin{align}
 \#\{\text{omitted coordinates}\}&\le5\tau^2v,
       \label{mr:matching-count}\\
 \|B-W\|_{\mathrm F}^2&\le8\tau^2v,
       \label{mr:matching-error}\\
 |B_{ij}-1|&\le\tau\quad\text{on every retained pair }\{i,j\}.
       \label{mr:matching-entry}
\end{align}
Every retained pair has $B_{ij}>0$.
\end{lemma}

\begin{proof}
Put
\[
 d_i=\sum_jB_{ij}^2,
 \qquad r_i=d_i^2-\sum_j B_{ij}^4\ge0.
\]
The diagonal entries of $B^2-I$ show that
$\sum_i(d_i-1)^2\le\varepsilon v$.  Nonnegativity also gives
\begin{align*}
 \sum_i r_i
 &=\sum_{j\ne k}\sum_iB_{ij}^2B_{ik}^2\\
 &\le\sum_{j\ne k}\left(\sum_iB_{ij}B_{ik}\right)^2
 \le\varepsilon v.
\end{align*}
Call row $i$ good when $|d_i-1|\le\tau$ and $r_i\le\tau^2$.
There are at most $2\tau^2v$ bad rows.  In a good row, let
$a_i=\max_j B_{ij}^2$.  Then
\[
 a_i\ge\frac{\sum_jB_{ij}^4}{d_i}
      =d_i-r_i/d_i,
 \qquad d_i-a_i\le\frac{\tau^2}{1-\tau}.
\]
Because $\tau\le1/8$, it follows that
\[
 a_i\ge\frac{1-2\tau}{1-\tau}\ge\frac67,
 \qquad d_i-a_i\le\frac87\tau^2,
 \qquad |\sqrt{a_i}-1|\le\tau.
\]
For the last inequality, use
$1-\tau-\tau^2/(1-\tau)\ge(1-\tau)^2$ and
$a_i\le1+\tau\le(1+\tau)^2$.
There is a unique entry attaining $a_i$.  If its other endpoint is
good, symmetry forces it to be the dominant entry there as well, since
its square is at least $6/7$ whereas every nondominant squared entry in
a good row is at most $8\tau^2/7$.  These good--good pairs form the
required matching.

For any set $S$ of rows, Cauchy--Schwarz gives
\begin{equation}\label{mr:row-mass}
 \sum_{i\in S}d_i
 \le |S|+\sqrt{|S|\varepsilon v}.
\end{equation}
Thus the total squared row mass at bad rows is at most
$(2+\sqrt2\tau)\tau^2v$.  Every good row whose dominant endpoint is
bad contributes at least $6/7$ to this mass, by symmetry.  The total
number of omitted coordinates is consequently at most
\[
 \left(2+\frac76(2+\sqrt2\tau)\right)\tau^2v
 \le5\tau^2v.
\]
For a retained row, its contribution to $\|B-W\|_{\mathrm F}^2$ is
at most $\tau^2+8\tau^2/7=15\tau^2/7$.  By
\eqref{mr:row-mass}, the omitted rows contribute at most
$(5+\sqrt5\tau)\tau^2v$.  Since $d\le v$ and $\tau\le1/8$, the
sum is less than $8\tau^2v$.  The remaining assertions have already
been proved.
\end{proof}

\begin{lemma}[A trace bound across the second normalization]
\label{mr:trace-cross}
For each $j$, let $H_j,Q_j$ be two lists of Lorentz vectors whose joint
Gram matrix has operator norm at most a fixed $K$.  Let $v_j\to\infty$.
Suppose that for every $\varepsilon>0$ the Gram matrix
$\langle Q_j,Q_j\rangle$ has $o(v_j)$ eigenvalues above
$2+\varepsilon$.  If
\[
 \limsup_j\frac{\Tr\langle H_j,H_j\rangle}{v_j}\le\eta,
\]
then
\begin{equation}\label{mr:trace-cross-bound}
 \limsup_j\frac{\|\langle H_j,Q_j\rangle\|_{\mathrm F}^2}{v_j}
 \le2\eta.
\end{equation}
The conclusion is independent of $K$.
\end{lemma}

\begin{proof}
Let $J_j$ be the joint Gram matrix.  Its negative part has rank at most
one and norm at most $K$.  Replacing $J_j$ by its positive part changes
any normalized trace by $O(K/v_j)$ and changes a Frobenius norm divided
by $\sqrt{v_j}$ by $O(K/\sqrt{v_j})$.  Write the resulting positive
semidefinite matrix as
\[
 P_j=\begin{pmatrix}A_j&C_j\\ C_j^{\mathsf T}&D_j\end{pmatrix}.
\]
A rank-one perturbation changes the number of eigenvalues above any
fixed threshold by at most one, by the variational characterization of
eigenvalues.  Thus the spectral projector $E_{\mathrm{hi}}$ of $D_j$
for $(2+\varepsilon,\infty)$ has rank $o(v_j)$.  Put
$E_{\mathrm{lo}}=I-E_{\mathrm{hi}}$.  Factor $P_j$ as a Euclidean
Gram matrix, say $A_j=U_j^{\mathsf T}U_j$,
$D_j=V_j^{\mathsf T}V_j$ and $C_j=U_j^{\mathsf T}V_j$.
Since $\|V_jE_{\mathrm{lo}}\|^2\le2+\varepsilon$,
\[
 \|C_jE_{\mathrm{lo}}\|_{\mathrm F}^2
 \le(2+\varepsilon)\Tr A_j.
\]
On the other spectral subspace,
$\|C_jE_{\mathrm{hi}}\|_{\mathrm F}^2\le K^2
\rank E_{\mathrm{hi}}=o(v_j)$.  Divide by $v_j$ and take the
limsup.  This first bounds $\|C_j\|_{\mathrm F}/\sqrt{v_j}$.
The original cross block differs by Frobenius norm at most $K$,
so it has the same squared normalized limsup bound, at most
$(2+\varepsilon)\eta$.  Finally let
$\varepsilon\downarrow0$.  In particular the hypotheses force
$\eta\ge0$, since the normalized trace of the original $H_j$ block
has nonnegative liminf.
\end{proof}

\subsection{The first normalization and the candidate groups}

\begin{proof}[Proof of \Cref{thm:matrix-obstruction}]
Use the constant $c=c(a)>0$ from \Cref{mt:weighted-bag}; explicitly we
may take
\begin{equation}\label{mr:constants}
 \begin{aligned}
 K(a)&=\left\lceil\frac5a\right\rceil
             \left(\left\lceil\log_2\frac5a\right\rceil+1\right),
 &c(a)&=\frac1{40K(a)},\\
 \delta_0(a)&=\min\left\{\frac1{3\cdot8^4},\frac{c(a)}{12},
                   \frac13\left(\frac{c(a)}{2340}\right)^4\right\},
 &\beta_0(a)&=\frac{\delta_0(a)}{10},\\
 m_0(a)&=\left\lceil\max\left\{\frac2a,
                                  \frac{10}{\delta_0(a)},2\right\}
                         \right\rceil.
 \end{aligned}
\end{equation}
Fix $a,\beta,m,G$ as in the theorem with these choices, and put
\begin{equation}\label{mr:error-parameters}
 q=\lceil\beta m\rceil,
 \qquad\delta=5(\beta+1/m),
 \qquad\tau=(3\delta)^{1/4}.
\end{equation}
These choices imply
\begin{equation}\label{mr:smallness}
 \delta\le\delta_0,
 \qquad\tau\le\frac18,
 \qquad\tau\le\frac{c}{2340},
 \qquad\frac qm\le\frac{\delta}{5}\le\frac c{60},
 \qquad am\ge2.
\end{equation}

Suppose the conclusion fails.  There is then an unbounded sequence of
integers $s$ and corresponding well-signed matrices $M_s$ with one
negative eigenvalue such that, writing $v=ms$,
\begin{equation}\label{mr:contradiction-rank}
 \rank M_s\le v/2+1.
\end{equation}
All subsequences below keep $a,\beta,m$ and the graph $G$ fixed, and
$s\to\infty$ is the only limiting operation.  An expression $o(v)$
has ratio to $v$ tending to zero in this limit; its constants may depend
on the fixed graph and subsequence.  The displayed constants in the
errors involving $\delta$ or $\tau$ are absolute.

Represent $M_s$ by Lorentz Gram vectors, one for each blow-up vertex.
An original vertex label $u\in V(G)$ specifies its $s$-column clique
block.  By passing to a subsequence, the set of labels whose clique
blocks have a negative eigenvalue can be fixed.  Two such labels cannot
be a hole: zero cross Gram would make their negative directions
orthogonal and give two negative eigenvalues of $M_s$.  These labels
therefore form a clique in $G$.  There are fewer than $2q$ of them,
since two disjoint $q$-element subsets of a larger clique would violate
condition (iii).  Discard these labels.

The other clique blocks are positive semidefinite.  They have strictly
negative off-diagonal entries; \Cref{mr:positive-block} shows that
deleting at most one coordinate makes each positive definite.  Work
with $s\ge2$, and apply \Cref{mt:finite-tree} to their positive spans.
For every hole between two labels the blocks are orthogonal, so their
tree subtrees intersect.  By \Cref{mt:single-bag}, there is a bag
missing fewer than $3q$ of the supplied single labels.  Keep the labels
in this bag and denote their set by $U$.  There are fewer than $5q$
missing original labels.  The retained coordinate dimension $d_s$
satisfies
\begin{equation}\label{mr:first-dimension}
 d_s\ge(m-5q)(s-1),
 \qquad \liminf_{s\to\infty}\frac{d_s}{v}\ge1-\delta.
\end{equation}
The tree and the chosen bag are fixed on a further subsequence.

Normalize each retained clique block by its positive inverse square
root, and denote the resulting column lists by $X_u$, $u\in U$.
Their diagonal Gram blocks are identities.  By
\Cref{mr:positive-block}, the normalizing coefficient matrices are
entrywise nonnegative with positive diagonal.  Thus cross blocks remain
zero on holes and strictly negative on base edges.  The full Gram matrix
of the $X_u$ is
\begin{equation}\label{mr:first-gram}
 A_s=I-B_s,
 \qquad B_s\ge0\ \text{entrywise},
 \qquad (B_s)_{uu}=0\ \text{as a whole diagonal block}.
\end{equation}
The bag bound from \Cref{mt:finite-tree,mt:flat-normalization} makes
$\|B_s\|$ uniformly bounded in $s$.  The rank is at most that of
$M_s$, and its negative index is at most one.  Consequently
\Cref{mr:spectral-support} gives
\[
 \limsup_{s\to\infty}\frac{\|B_s^2-I\|_{\mathrm F}^2}{v}
 \le2\delta.
\]
Since $\delta>0$ is fixed, the left-hand quantity before taking a
limsup is at most $3\delta=\tau^4$ for all sufficiently large $s$.
Apply \Cref{mr:matching} to obtain a coordinate matching with $p_s^0$
pairs and adjacency matrix $W_s$.  Then
\begin{equation}\label{mr:first-matching}
 \begin{split}
 \liminf_s\frac{p_s^0}{v}
     &\ge\frac12-\frac\delta2-\frac52\tau^2,\\
 \|B_s-W_s\|_{\mathrm F}^2&\le8\tau^2v,
 \qquad |(B_s)_{ij}-1|\le\tau
          \quad\text{on retained pairs}.
 \end{split}
\end{equation}
Every pair connects two different adjacent base labels, because
$B_s$ vanishes within each original block and on base holes.

Group the pairs by their unordered label pair $e=\{u,w\}\in E(G)$.
If group $e$ contains $k_e(s)$ pairs, pass to a subsequence on which
every $k_e(s)/s$ converges, and put
\begin{equation}\label{mr:weights}
 w_e=\lim_s k_e(s)/s,
 \qquad \sum_{e\ni u}w_e\le1,
 \qquad \sum_e w_e\le m/2.
\end{equation}
There are only finitely many types, so we may also fix whether each
group is empty or its endpoint Gram block is positive definite,
singular positive semidefinite, or has a negative eigenvalue.

Groups in the last category will not be candidates for the second
normalization.  Their types are pairwise touching: a conflict between
two of them would make their cross Gram zero and give two orthogonal
negative directions.  Choose a maximal ordinary matching of these
types in the base graph, of size $k$.  If $k\ge m/20$, putting each edge
in its own class contradicts (i), since $am\ge1$.  Thus $k<m/20$.
The $2k$ endpoints of this maximal matching cover all the types under
consideration.  The load bound in \eqref{mr:weights} therefore gives
\begin{equation}\label{mr:indefinite-weight}
 \sum_{e:\,\text{group has a negative eigenvalue}}w_e
       \le2k<m/10.
\end{equation}
Their coordinate pairs are retained in the first matching.

For every group of positive weight with positive semidefinite endpoint
Gram, that Gram has the form
\[
 \begin{pmatrix}I&-C\\-C^{\mathsf T}&I\end{pmatrix},
 \qquad C\ \text{entrywise strictly positive}.
\]
The graph of negative entries is connected.  If the group is singular,
delete one whole matching pair.  By \Cref{mr:positive-block} its
remaining Gram block is positive definite.  Positive weight ensures
that the group is eventually nonempty after this deletion.  There are
at most $\binom m2$ groups, so these deletions cost $O(m^2)=o(v)$
coordinates.  Call the resulting positive definite groups of positive
weight the \emph{candidate groups}.  Groups of zero limiting weight
need not be candidates and keep their pairs in the first matching.
By \eqref{mr:first-matching}, \eqref{mr:indefinite-weight} and
$\tau\le1/8$, the total candidate weight satisfies
\begin{equation}\label{mr:candidate-weight}
 \sum_{e\text{ candidate}}w_e
 \ge m\left(\frac25-\frac\delta2-\frac52\tau^2\right)
 >\frac{3m}{10}.
\end{equation}

\subsection{A common bag for both normalizations}

Apply \Cref{mt:finite-tree} again, this time to the single-label
positive spans for all $u\in U$ together with the positive spans of all
candidate groups.  These spans are allowed to overlap.  The required
subtree contacts hold in each of the following cases: two single labels
form a hole; two group types conflict; or a group type and a single
label conflict.  In each case all cross Gram entries vanish, so the
positive spans are orthogonal and their flats intersect.

The candidate weights satisfy \Cref{mt:weighted-bag}, so a bag has
group weight at least $cm$.  Apply \Cref{mt:common-bag} to the supplied
single labels $U$ and these groups, using $q\le cm/60$.  It yields a
bag containing all but fewer than $5q$ labels of $U$, and containing
candidate group weight at least $cm/2$.  The common-bag proof uses all
root labels absent after its threshold edge, so it applies equally
when the losses occur one label at a time.

Discard from the remaining coordinate matching every pair incident to one of the
missing labels of $U$.  This removes whole label-pair groups and at most
$5qs$ pairs.  Denote the surviving total number of pairs by $p=p_s$.
Let $\mathcal T$ be the candidate groups in the common bag whose
endpoints both remain, and let $t=t_s$ be their total number of pairs.
Since the load at a missing label is at most one, the limiting weight
lost from the bag is at most $5q$.  Thus
\begin{align}
 \liminf_s\frac pv
 &\ge\frac12-\frac\delta2-\frac52\tau^2-\frac{5q}{m}
 \ge\frac12-3\tau^2,
       \label{mr:final-pair-count}\\
 \lim_s\frac tv
 &\ge\frac c2-\frac{5q}{m}\ge\frac c3.
       \label{mr:extra-pair-count}
\end{align}
For the last inequality in \eqref{mr:final-pair-count}, use
$5q/m\le\delta=\tau^4/3$ and $\tau\le1$.  All group counts have
limits after the earlier subselection, so $t/v$ does as well.

Write $X$ for the list of the $2p$ surviving matched endpoint columns,
still in the first normalization. For $e\in\mathcal T$, let $k_e$
denote its remaining number of pairs and write $X_e$ for its
$2k_e$-column sublist, $D_e=\langle X_e,X_e\rangle$, and
\begin{equation}\label{mr:second-normalization}
 Q_e=X_eD_e^{-1/2},
 \qquad Q=(Q_e)_{e\in\mathcal T}.
\end{equation}
The common bag gives a uniform operator bound for the joint Gram of
the first-normalized single blocks and all $Q_e$.  Passing to the
sublist $X$ does not increase this bound.  This point controls $Q_e$
even when some eigenvalues of $D_e$ tend to zero and the coefficients
in $D_e^{-1/2}$ become large.

\begin{figure}[tb]
\centering
\begin{tikzpicture}[
  box/.style={draw,rounded corners,align=center,text width=4.2cm,
              minimum height=1.85cm,inner sep=5pt},
  >=Latex,font=\small]
\node[box] (first) {First normalization\\
  $\langle X_u,X_u\rangle=I$\\coordinate pairs};
\node[box,right=1.1cm of first] (sumdiff) {Pair differences $R$\\
  $\langle R,R\rangle\approx2I$\\pair sums $H$\\have small trace};
\node[box,below=1.15cm of first] (second) {Second normalization\\
  $\langle Q_e,Q_e\rangle=I$\\within candidate groups};
\node[box,right=1.1cm of second] (extra) {Exact complements $Q_eL_e$\\
  $k_e$ extra directions\\orthogonal to $R_e$};
\draw[->] (first)--(sumdiff);
\draw[->] (first)--node[left,align=right]{one common\\tree bag}(second);
\draw[->] (second)--(extra);
\draw[->] (sumdiff)--node[right,align=left]{small cross Gram\\via $H$}(extra);
\end{tikzpicture}
\caption{The two normalizations.  The common bag bounds all indicated
Gram matrices.  The trace estimate for $H$, followed by a spectral cutoff
at $2$, controls the interaction with the second normalization using an
absolute constant.}
\label{mr:two-scales}
\end{figure}

\subsection{The additional positive directions}

For each surviving pair $j=\{i,i'\}$ choose an order of its endpoints
and form
\[
 R_j=\frac{X_i-X_{i'}}{\sqrt2},
 \qquad H_j=\frac{X_i+X_{i'}}{\sqrt2}.
\]
The coefficient maps producing either list from $X$ are Euclidean
isometries.  The restriction of $I-W_s$ to retained pair endpoints has
Gram $2I$ on the difference coefficients.  Hence
\eqref{mr:matching-error} and contraction of Frobenius norm give
\begin{equation}\label{mr:R-bound}
 \|\langle R,R\rangle-2I_p\|_{\mathrm F}^2\le8\tau^2v.
\end{equation}
Also $\langle H_j,H_j\rangle=1-(B_s)_{ii'}$, so
\begin{equation}\label{mr:H-trace}
 |\Tr\langle H,H\rangle|\le p\tau\le\frac{\tau v}{2}.
\end{equation}
These statements persist after all the pair deletions: they are
restrictions of the original matching estimates.

The coefficients in each $D_e^{-1/2}$ are entrywise nonnegative.
Distinct groups have disjoint endpoint coordinates, and all cross Gram
entries between such coordinates are nonpositive.  Coordinates with
the same base label have zero interaction, since the first diagonal
blocks are identities.  Therefore
\begin{equation}\label{mr:Q-sign}
 \langle Q,Q\rangle=I-C_s,
 \qquad C_s\ge0\ \text{entrywise},
 \qquad (C_s)_{ee}=0\ \text{as a whole group block}.
\end{equation}
It is bounded in operator norm and has at most one negative eigenvalue.
By \Cref{mr:spectral-support}, its limiting spectral support lies in
$[0,2]$: more precisely, for every fixed $\varepsilon>0$ there are
$o(v)$ eigenvalues outside $[-\varepsilon,2+\varepsilon]$.
The joint Gram of $(H,Q)$ is bounded by the common bag.  Applying
\Cref{mr:trace-cross} to \eqref{mr:H-trace} yields the absolute estimate
\begin{equation}\label{mr:HQ-bound}
 \limsup_s\frac{\|\langle H,Q\rangle\|_{\mathrm F}^2}{v}
 \le\tau.
\end{equation}
In particular, no factor depending on the norm of the common bag
appears here.

Fix a candidate group $e\in\mathcal T$. Let $U_e$ be the
$2k_e\times k_e$ Euclidean isometry whose columns are the coefficient
vectors of its pair differences. Then
\[
 R_e=X_eU_e=Q_e T_e,
 \qquad T_e=D_e^{1/2}U_e.
\]
Because $D_e$ is positive definite, $T_e$ has rank exactly $k_e$.
Choose a $2k_e\times k_e$ matrix $L_e$ whose columns are an orthonormal
basis of $(\im T_e)^\perp$.  Thus
\begin{equation}\label{mr:exact-complement}
 L_e^{\mathsf T}L_e=I_{k_e},
 \qquad\langle R_e,Q_eL_e\rangle=T_e^{\mathsf T}L_e=0.
\end{equation}
This uses the full span of all the group's original difference
directions; no small singular direction is discarded.  Write
$L=\diag(L_e)_{e\in\mathcal T}$ and $Y=QL$.  The list $Y$ has
exactly $t$ columns.

For a pair $j$ outside group $e$, each of the two numbers
$x=\langle X_i,(Q_e)_k\rangle$ and
$y=\langle X_{i'},(Q_e)_k\rangle$ is nonpositive.  Consequently
\begin{equation}\label{mr:sign-domination}
 |\langle R_j,(Q_e)_k\rangle|
    =\frac{|x-y|}{\sqrt2}
    \le-\frac{x+y}{\sqrt2}
    =-\langle H_j,(Q_e)_k\rangle.
\end{equation}
This includes pairs in noncandidate groups.  Let $Z$ be
$\langle R,Q\rangle$ with the entries between each $R_e$ and its
own $Q_e$ replaced by zero.  Then
$\|Z\|_{\mathrm F}^2\le\|\langle H,Q\rangle\|_{\mathrm F}^2$,
and \eqref{mr:exact-complement} gives
$\langle R,Y\rangle=ZL$.  Since $L$ is an isometry, Frobenius
contraction and \eqref{mr:HQ-bound} imply
\begin{equation}\label{mr:RY-bound}
 \limsup_s\frac{\|\langle R,Y\rangle\|_{\mathrm F}^2}{v}
 \le\tau.
\end{equation}

Put $E_s=\langle Y,Y\rangle=L^{\mathsf T}\langle Q,Q\rangle L$.
Every diagonal group block of $E_s$ is an identity, so
\begin{equation}\label{mr:Y-trace}
 \Tr E_s=t.
\end{equation}
This is a Euclidean coefficient compression.  The variational
characterization of eigenvalues shows that the number of eigenvalues
of $E_s$ above $2+\varepsilon$ is at most that for
$\langle Q,Q\rangle$, hence is $o(v)$.  Its norm is bounded.
If $k_s$ denotes its number of eigenvalues greater than $1/4$, use
the cutoff $2+\varepsilon=5/2$ to obtain
\[
 t=\Tr E_s\le\frac{t-k_s}{4}+\frac52k_s+o(v).
\]
Possible negative eigenvalues only decrease the left-hand trace
contribution in this upper estimate.  Rearranging gives
\begin{equation}\label{mr:Y-positive-count}
 k_s\ge t/3-o(v).
\end{equation}

Finally consider the joint Gram matrix of $(R,Y)$.  By
\eqref{mr:R-bound}, the number of eigenvalues of
$\langle R,R\rangle$ at most $1/4$ is at most
\[
 \frac{16}{49}\|\langle R,R\rangle-2I_p\|_{\mathrm F}^2
 \le\frac{128}{49}\tau^2v.
\]
Restrict the two diagonal blocks to their spectral spaces above
$1/4$.  This leaves at least
$p-(128/49)\tau^2v$ dimensions on the $R$ side and
$t/3-o(v)$ on the $Y$ side.  The restricted cross matrix has squared
Frobenius norm at most $\|\langle R,Y\rangle\|_{\mathrm F}^2$.
Delete from the first side its left singular directions with singular
values greater than $1/8$.  The number deleted is at most
$64\|\langle R,Y\rangle\|_{\mathrm F}^2$.  The remaining diagonal
blocks are greater than $I/4$, and the cross operator norm is at most
$1/8$.  For their coefficient vectors $x,y$, the joint quadratic form
is bounded below by
\[
 \frac14(\|x\|^2+\|y\|^2)-\frac14\|x\|\|y\|
 \ge\frac18(\|x\|^2+\|y\|^2),
\]
so this remaining space is positive definite.  Using
\eqref{mr:RY-bound}, the positive index of the joint Gram is at least
\begin{equation}\label{mr:rank-gain}
 p+\frac t3-\frac{128}{49}\tau^2v-64\tau v-o(v).
\end{equation}
Every column used here is a linear combination of the original Lorentz
Gram columns for $M_s$.  The joint Gram is therefore a coefficient
pullback of $M_s$, and its positive index is at most $\rank M_s$.
This argument does not require the combined list $(R,Y)$ to be
linearly independent.

Combine \eqref{mr:final-pair-count}, \eqref{mr:extra-pair-count} and
\eqref{mr:rank-gain}.  Since $\tau\le1/8$,
\begin{align*}
 \liminf_s\frac{\rank M_s}{v}
 &\ge\frac12+\frac c9-64\tau-\frac{275}{49}\tau^2\\
 &\ge\frac12+\frac c9-65\tau
 \ge\frac12+\frac c{12},
\end{align*}
where the last inequality uses $\tau\le c/2340$.
This contradicts \eqref{mr:contradiction-rank}, whose normalized
right-hand side tends to $1/2$.  Hence only finitely many $s$ can
violate the stated rank bound.  The positive gap also shows directly
why the literal additive $+1$ is harmless for all sufficiently large
$s$.

Every matrix in the definition of $\mu(G[K_s])$ is among the matrices
just considered.  Rank--nullity gives
$\dim\ker M<ms/2-1$, and taking the maximum yields
$\mu(G[K_s])+1<ms/2$ as asserted.
\end{proof}

For completeness, the matrix class defining $\mu(H)$ is nonempty for
every finite nonempty graph $H$.  Start from a diagonal matrix with one diagonal
entry $-1$ and all others $1$, and put the same sufficiently small
negative number in each edge position.  The perturbation has operator
norm less than $1/2$ when its size is sufficiently small, so the matrix
is nonsingular with exactly one negative eigenvalue.  Nonsingularity
makes the Strong Arnold Property immediate.

If the base graph additionally satisfies $\alpha(G)\le2$, then
$\alpha(G[K_s])\le2$: an independent set uses at most one vertex of
each clique, and its base labels are independent.  Thus
$\chi(G[K_s])\ge\lceil ms/2\rceil>\mu(G[K_s])+1$.  Together with
\Cref{thm:base-graphs}, the theorem above proves \Cref{thm:main}.
One may also recover the chromatic number exactly as
$\chi(G[K_s])=ms-\nu(\overline{G[K_s]})$, where $\nu$ is the maximum
matching size: every color class has one or two vertices, and the
two-vertex classes are precisely disjoint hole pairs.
